\documentclass{report}
\usepackage{amsmath,amssymb}
\setlength{\parindent}{0mm}
\setlength{\parskip}{1em}
\begin{document}
\begin{center}
\rule{6in}{1pt} \
{\large Giovanni De Luca \\
{\bf GLOBAL DISCRETE EMPIRICAL INTERPOLATION METHOD FOR NONLINEAR MODEL ORDER REDUCTION}}

Avenue des Jeux Olympiques \\ 2
\\
{\tt g.de.luca@tue.nl}\\
Michiel E. Hochstenbach\end{center}

In the context of model order reduction (MOR) for nonlinear problems, we
propose a variant of discrete empirical interpolation method (DEIM),
called global DEIM (GDEIM). Since DEIM is a nearly-optimal method, i.e.,
the error of the reduction depends on the number and location of the
selected nonlinearities to evaluate, we show that, by employing an
alternative procedure to select the points, we sometimes obtain a better
accuracy of the reduction with GDEIM for the same number of points as for
DEIM. This is useful especially when one has restrictions on the
dimension of the reduced order model (ROM) to simulate or, the other way
around, when a better accuracy can be reached for a speci ed reduction
order. To compare the accuracy level between GDEIM and DEIM before
proceeding with the reduction, i.e., when the projection matrix is
formed, we make use of a heuristic error estimator cheap to evaluate,
which in most of our experiments is able to determine which of the two
methods is best. This can speedup simulations of the obtained ROM, when
inputs and/or parameters are swept for analysis. We demonstrate the
usefulness of GDEIM to nonlinear, parametric, analytic functions and
provide a comparison with DEIM by using our error estimator.
The second contribution of this work is a mathematical formulation which
extends (G)DEIM to general nonlinear function, evaluated
non-componentwise. We provide accuracy and speedup results from using our
extended formulation to an example for time-domain circuit simulation,
described by a set of differential algebraic equations. This may be a
promising direction to speedup analysis of more complex and large
circuits, where it is common to have millions of unknowns in the
vector-solution.


\end{document}
